\section{Information bounds for separated regions}
\label{sec:information}

Throughout this section, distances in the plane are measured in the sup
norm.  The distance used in a Hamiltonian propagation estimate is the
nearest-neighbor graph distance on the physical square; every graph ball
is contained in the corresponding ambient sup-norm ball.  As in Section~\ref{sec:preliminaries}, $F^{+r}$ denotes the ambient
radius-$r$ enlargement of a plane set $F$.  Its physical sample is
$\Lambda\cap F^{+r}$.  Ambient rectangles and squares include their
unoccupied locations as well.  Constants
may depend on $q,J,\Delta$ and on explicitly fixed geometric parameters,
but not on $L$ or on the Hamiltonian.  A lower threshold on $L$, or on a
geometric scale, always has the same uniformity.

The purpose of this section is to prove that information is negligible
between regions whose separation decreases linearly toward a fixed
number of points, provided the separation is never smaller than a
sufficiently large multiple of $\log L$.  We first extract a coarse
collar estimate from the companion results, then prove a conditional
information bound for one cell.  A separate consequence of the global
gap, using Gaussian spectral filtering as in \cite{Hastings2006Local}, reduces an information cost to an arbitrarily small inverse
power of $L$ after shrinking the regions by a width proportional to that
cost plus $\log L$.  Applying that reduction after each annular step gives the
required logarithmic minimum separation.

\subsection{The coarse collar interface}

We specify the additional companion statements used here, beyond the
entropy area law in Proposition~\ref{prop:area-law}.  The constants
$C_{\rm tpl}$ and $D_0$ below are fixed template constants from the companion.
In the terminology
of \cite[Sections~9 and~10]{RepoAreaLaw}, a template is a nonempty ambient
lattice set
\[
 T=\bigcup_{a=1}^{N_0}(P_a\cap\mathbb Z^2),
 \qquad
 n\ge C_{\rm tpl}N_0(s_0+1),
 \qquad \diam_\infty T\le n,
\]
where the pieces are closed rectangles or triangles, have sup diameter
at most $s_0$, and have sides of slopes $0,\infty,1,-1$.  The companion also
fixes a physical cut $A\subseteq\Lambda$ and its set $Z_A$ of
crossing-edge endpoints.  Its template-clearance assumption is
$\dist_\infty(T,Z_A)>4D_0s_0$.  Put
\[
 \varepsilon_{\rm col}=10^{-5},\qquad
 \ell=\lfloor n^{1-\varepsilon_{\rm col}}\rfloor,
 \qquad T_\ell=T+([-\ell,\ell]^2\cap\mathbb Z^2).
\]
The Proposition ``A sublinear collar bound'' in
\cite[Section~9]{RepoAreaLaw} states that, when $\ell\le s_0$ and the
scale is sufficiently large,
\begin{equation}
 \begin{split}
 X&=A\cap T,\qquad S_0=A\cap T_\ell,\\
 I_\Omega(X:\Lambda\setminus S_0)&\le Cn^{1-\varepsilon_{\rm col}}.
 \end{split}
 \label{eq:info-repository-collar}
\end{equation}
It also supplies $|S_0|\le Cn^2$ and $Cn$ upper bounds for the numbers
of physical crossing edges of all three cuts $X,S_0,S_0\setminus X$.
All constants and thresholds are uniform in the domain, cut, and
admissible template.

The Proposition ``Amplification'' in \cite[Section~10]{RepoAreaLaw}
uses precisely these size, three-cut, and information hypotheses, with
any fixed common upper-bound constant.  Its fixed exponents are
\[
 \beta_0=1-2\cdot10^{-6},\qquad
 \alpha=\frac{k}{k+1}>1-10^{-6},
\]
where $k$ is a sufficiently large fixed integer.  For every fixed integer
$M\ge1$, it gives a graph radius $r$ such that
\begin{equation}
 \ell+2r=o(n^{\beta_0}),\qquad
 I_\Omega(X:E)\le C_Mn^{-M}
 \quad\text{whenever }E\subseteq\Lambda\setminus N_{2r}(S_0).
 \label{eq:info-repository-amplification}
\end{equation}
Here $N_{2r}$ is a physical graph neighborhood.  The positive system
required by that Proposition is furnished by the Proposition ``Positive
replacement of the Hamiltonian'' in \cite[Section~4]{RepoAreaLaw}, for
each prescribed integer $k$.  That replacement assumes the original
full-system eigenvector and gap inequality.  Its constants may depend
on $k$, which is fixed here.  It does not assume a gap for any restricted
Hamiltonian.  Our on-site and nearest-neighbor Hamiltonian has
induced-graph range one and one bounded term per designated support, so
it satisfies the bounded-interaction hypotheses of all these results.

\begin{proposition}[Coarse information outside a sublinear collar]
\label{prop:coarse}
Fix an integer $N_0\ge1$ and a constant $A_0\ge1$.  There are
$\beta<1$, $C<\infty$, and $u_0<\infty$ such that the following holds.
Let $u\ge u_0$ and let $T$ be the full ambient lattice sample of a union
of at most $N_0$ axis-parallel rectangles, with
$\diam_\infty T\le A_0u$.  For every physical set $E\subseteq\Lambda$
satisfying
\[
 \dist_\infty(E,T)>u^\beta,
\]
one has
\[
 I_\Omega(\Lambda\cap T:E)\le C.
\]
The same statement holds for any subset of $\Lambda\cap T$ in its first
argument, for rectangles with consistently assigned open or closed
sides, and after a common translation of the lattice coordinates.
\end{proposition}

\begin{proof}
An empty ambient sample gives the assertion immediately.  Otherwise
omit the empty pieces, choose an integer $s_0$ comparable to $u$ and
large enough to bound every piece diameter, and then choose an integer
$n$ comparable to a sufficiently large fixed multiple of $u$.  These
choices ensure both template inequalities.  The prescribed cut is
$A=\Lambda$.  Consequently $Z_A=\varnothing$, so all template-clearance
conditions are vacuous.  Since the ratio $n/s_0$ is fixed,
\[
 \lfloor n^{1-\varepsilon_{\rm col}}\rfloor\le s_0
\]
for all sufficiently large $u$.  Equation~\eqref{eq:info-repository-collar}
therefore supplies every hypothesis of
\eqref{eq:info-repository-amplification}; fix its common constant once
and apply it, for example, with $M=1$.

Choose $\beta$ strictly between $\beta_0$ and $1$.  Because $n$ is a
fixed multiple of $u$,
\[
 \ell+2r<u^\beta
\]
for sufficiently large $u$.  A graph neighborhood of $S_0$ of radius
$2r$ is contained in the physical sample of the ambient enlargement of
$T$ by $\ell+2r$.  Hence every $E$ in the statement is an allowable
remote set in \eqref{eq:info-repository-amplification}.  The bound
$C_1n^{-1}$ is in particular uniformly bounded.  Discarding part of the
first system cannot increase mutual information.

On the lattice, a rectangle with some open sides has the same sample
as a closed rectangle with suitably shifted sides; the shifts can be
arbitrarily small unless an endpoint lies on a lattice line, in which
case they can still be chosen without changing the selected sites.
The diameter bounds change by at most a fixed additive amount.  A
common coordinate translation can be undone before using the companion
statements.  Increasing the threshold and constants proves these
versions too.
\end{proof}

\subsection{A conditional information bound for a cell}

The next construction adapts the two-family ordering in
\cite[Section~11]{RepoAreaLaw} to a cell with arbitrary exterior conditioning.
We prove the construction and its conditioning order here.
We use ambient colors in the next construction.  Missing physical sites
are simply ignored when taking entropies.  Boundaries of all chambers
are assigned by one consistent infinitesimal generic displacement;
equivalently, one works with open chambers, uses their closures for
distance estimates, and then samples with a common tie convention.
Only finitely many chamber boundaries occur at a fixed instance.  No
color is assigned exclusively to a lower-dimensional face.

\begin{lemma}[Conditional information of one cell]
\label{lem:cell-cmi}
Fix $c>0$.  Tile the plane by an axis-parallel grid of squares of side
$s\ge1$.  Let $G$ be the physical sample of one cell, let $B$ be the
physical sample of any union of other cells, and let
$C\subseteq\Lambda\setminus(B\cup G)$ satisfy
\[
 \dist_\infty(C,\overline{G}_{\rm amb})\ge cs,
\]
where $\overline{G}_{\rm amb}$ denotes the full closed ambient cell.
There are fixed constants $C_1,K<\infty$ such that
\begin{equation}
 I_\Omega(G:C\mid B)\le C_1(1+\log s)^K.
 \label{eq:info-cell-cmi}
\end{equation}
The constants are independent of which exterior cells comprise $B$
and of the shape of $C$ away from the cell.
\end{lemma}

\begin{proof}
Outside the chosen cell, color all cells of $B$ by forbidden color $0$
and all other cells by forbidden color $1$.  In particular this colors
also the physical sites outside $B,G,C$.  The interior of the chosen
cell is initially allowed.  We construct two ordered tile families
$\mathcal F_0,\mathcal F_1$, together with a set of exceptional sites.
The order in a family is its order of creation.  Whenever an allowed
region is assigned to family $i$, give it forbidden color $i$ for all
subsequent stages.  A tile $U$ in family $i$ will satisfy
\begin{equation}
 I_\Omega\bigl(U:C\cup F_i\cup
       \{\text{earlier tiles of family }i\}\bigr)\le C_2,
 \label{eq:info-tile-cost}
\end{equation}
where $F_i$ is the physical exterior of the initial cell of forbidden
color $i$.  Tile notation in an entropy denotes its physical sample.

Assume for the moment that such a partition has been constructed.  Its
information cost bounds the desired conditional information as follows.
Apply the chain rule
for $I(G:C\mid B)$ by listing first $\mathcal F_0$ in creation order,
then all exceptional sites, then $\mathcal F_1$ in reverse creation
order.  For a tile $U$ the chain-rule term is $I(U:C\mid Z)$, where
$Z$ is $B$ together with previously listed sites.  In family $0$,
\[
 I(U:C\mid Z)=I(U:CZ)-I(U:Z)\le I(U:CZ),
\]
and $Z$ consists of forbidden-$0$ exterior sites and earlier
family-$0$ tiles.  Equation~\eqref{eq:info-tile-cost} applies.
For a family-$1$ tile, purity on the whole physical system gives
\[
 I(U:C\mid Z)=I\bigl(U:C\mid(UCZ)^c\bigr).
\]
The complementary conditioning set contains only forbidden-$1$
exterior sites and as-yet unprocessed family-$1$ tiles.  Because this
family is processed in reverse order, those tiles were created earlier
than $U$.  Bounding the last conditional mutual information by the
ordinary mutual information with $C\cup(UCZ)^c$ again invokes
\eqref{eq:info-tile-cost}.  Each exceptional one-site term is at most
$2\log q$, by the conditional entropy dimension bound.  Thus the desired
bound follows once the number of tiles and exceptional sites is at most
$C(1+\log s)^K$.  This explains both the two exterior
colors and the opposite orders in which the two families will be read.

We now construct that partition by an ambient recursion.  Fix
\[
 \beta<\gamma<\gamma'<1,
\]
with $\beta$ from Proposition~\ref{prop:coarse} for bounded rectangular
unions.  At a given generation, all allowed residual material lies in
active axis-parallel squares of a common side $u$.  Distinct active
squares have mutual distance greater than $10u$.  In the ambient
radius-$u$ padding of each active square, the three labels (allowed,
forbidden $0$, forbidden $1$) are constant on chambers of boundedly
many horizontal and vertical lines.  The square sides are included,
and distinct parallel lines are at least $u/8$ apart.  Initially there
is one active square, the original grid cell, with $u=s$.  Its padded
window meets only boundedly many grid lines, which verifies the
invariant regardless of the exterior color choices.

Choose a sufficiently large fixed stopping threshold $u_*$.  At a
nonterminal square put $t_0=u^\gamma$.  Increase $u_*$ so that all
inequalities below hold for $u\ge u_*$.  Around each represented line
intersection within distance $10t_0$ of the square, reserve a closed
square of radius $4t_0$ for the next generation.  Reserving a hole does
not change a previously fixed color; only still-allowed material is
reserved.  The hole squares are disjoint.  Every allowed point of the
active square outside these holes can see at most one forbidden color
within distance $t_0$.  Indeed it cannot see two parallel represented
lines, because $2t_0<u/8$.  If it sees both orientations, their
intersection is within distance $t_0$ of the point and belongs to a
reserved hole, a contradiction.  Thus its neighborhood sees at most
one line.  The chamber on its own side is allowed, and the chamber on
the other side has at most one forbidden color.

Using the colors as they were at the start of this generation, assign
to family $1$ every remaining allowed point within distance $t_0$ of
forbidden color $0$.  Assign all other remaining allowed points to
family $0$.  The new family-$0$ tile is at least $t_0$ from the old
forbidden-$0$ set.  The new family-$1$ tile is at least $t_0$ from the
old forbidden-$1$ set, by the preceding one-color observation.  There
is at most one tile of each family per active square.  Tiles of the
same family in different active squares are farther apart still, so
they may be ordered arbitrarily within the generation.  The distance
tests in the sup norm are unions of axis-parallel box enlargements;
after also deleting the bounded list of holes, each new tile is a union
of boundedly many rectangular chambers of total diameter $O(u)$.

The ambient template of such a tile is at distance at least $t_0$ from
the physical old forbidden set of its own family.  It is also at
distance at least $cs$ from $C$, since the tile lies in the original
cell.  Since $u\le s$ and $\gamma>\beta$, both distances exceed the
collar required by Proposition~\ref{prop:coarse} after increasing $u_*$.
This proves \eqref{eq:info-tile-cost} for every new tile.  In particular,
it bounds the information with the union of all earlier tiles of that
family, not just with any one earlier tile.

The new active squares are the reserved holes that contain allowed
material.  Their side is
\[
 u'=8t_0.
\]
Distinct represented intersections are separated by at least $u/8$.
Thus, for sufficiently large $u_*$, the new squares are at mutual
distance greater than $10u'$.  Squares descending from different
parents inherit even larger separation.  To verify the chamber
invariant, consider a child's radius-$u'$ padded window, of radius
$12t_0$ about its center.  The old line spacing implies that the only
old horizontal and vertical lines in this window are the lines through
its center.  New boundaries in the window occur only at offsets
\[
 0,\quad \pm t_0,\quad \pm4t_0
\]
from those two lines: these account for the distance test, the reserved
hole, and any clipping by the old active-square sides.  Other holes
and old parallel lines are outside the window.  Consequently the new
parallel-line spacing is at least $t_0=u'/8$, and the number of lines
is uniformly bounded.  This proves the invariant at the next generation.
The same verification covers a reserved hole centered on or slightly
outside a parent side; old fixed colors are retained and only the
allowed intersection is processed.

At scale below $u_*$, stop and make every remaining physical site
exceptional.  There are boundedly many children per active square and
boundedly many sites in each terminal square.  Choose $u_*$ also so that
\[
 8u^\gamma\le u^{\gamma'}\qquad(u\ge u_*).
\]
The number of generations is $O(1+\log\log(s+e))$.  Bounded branching
therefore bounds the total number of tiles and exceptional sites by
$C(1+\log s)^K$ for some fixed $C,K$.  If $s<u_*$, the direct bound
$I(G:C\mid B)\le2|G|\log q$ already has this form.  The boundary
convention stated before the Lemma preserves all separation estimates
on closures.  Each sampled tile remains a sampled rectangular union,
so it is covered by Proposition~\ref{prop:coarse}.  The entropy reduction
at the start of the proof now gives \eqref{eq:info-cell-cmi}.
\end{proof}

\subsection{Reducing a collar information cost with the global gap}

We use Lemma~\ref{lem:fidelity} in the following form.  For density
operators with $\supp\rho\subseteq\supp\sigma$,
\begin{equation}
 \norm{\sqrt\rho\sqrt\sigma}_1
 \ge \exp\left(-\tfrac12\Tr\rho(\log\rho-\log\sigma)\right),
 \label{eq:info-relative-entropy-overlap}
\end{equation}
and, after enlarging the purifying space if necessary, an isometry on
that space realizes the corresponding overlap with a prescribed
purification.  In particular the original purifying buffer need not
be minimal.  All enlarged factors will be removed before a Hamiltonian
locality estimate is applied.

\begin{lemma}[Information reset]
\label{lem:reset}
For every fixed $p>0$ there is $C_p<\infty$ with the following
property for all sufficiently large $L$.  Let $R,C\subseteq\Lambda$
be disjoint, let $R_0\subseteq R$ and $C_0\subseteq C$, and suppose
\[
 I_\Omega(R:C)\le b,\qquad b\ge0.
\]
If
\begin{equation}
 \dist_\infty(R_0,\Lambda\setminus R)>w,\qquad
 \dist_\infty(C_0,\Lambda\setminus C)>w,
 \qquad w=C_p(1+b+\log L),
 \label{eq:info-reset-width}
\end{equation}
then
\[
 I_\Omega(R_0:C_0)\le L^{-p}.
\]
Only the original full-system gap is assumed.
\end{lemma}

\begin{proof}
If either core is empty there is nothing to prove.  Use two replicas
and put
\[
 \Psi=\Omega^{\otimes2},\qquad F=F_R,\qquad \Phi=F\Psi,
\]
where $F_R$ exchanges the two replicas on $R$.  Our aim is to construct
a unitary $U$, supported outside $R_0\cup C_0$ in both replicas, for
which $U\Psi$ is close to $\Phi$.  Such a unitary leaves the joint core marginal in one replica
unchanged.  In $\Phi$, however, $R_0$ and $C_0$ come from different
original replicas, so this marginal is $\rho_{R_0}\otimes\rho_{C_0}$.
Closeness of the two vectors will therefore force small information
between the cores.

We first obtain a controlled overlap using an operator on the physical
buffer $Q=(R\cup C)^c$ of the two replicas.  Spectral filtering and
localization will turn it into a two-sided approximation to the desired
map; its polar factor will supply $U$ on the same physical support.

The relative entropy of $\rho_{RC}$ from
$\rho_R\otimes\rho_C$ is $I_\Omega(R:C)$.  The required support
inclusion holds: positivity implies that a joint state is supported
on the tensor product of the supports of its marginals.  For example,
the expectation of the positive joint state on the kernel projector
of either marginal is zero, so that projector annihilates the joint
state.  Equation~\eqref{eq:info-relative-entropy-overlap} and the
purification formula therefore give an isometry
\[
 V:\mathcal H_Q\longrightarrow\mathcal H_{b_R}\otimes\mathcal H_{b_C}
\]
and normalized vectors $a_{Rb_R}$, $a'_{Cb_C}$ such that
\begin{equation}
 \bigl|\langle a\otimes a',(\id_{RC}\otimes V)\Omega\rangle\bigr|
       \ge e^{-b/2}.
 \label{eq:info-split-overlap}
\end{equation}
The factors $b_R,b_C$ may first be enlarged by blank ancillary factors
so that the product purification and an isometric embedding of the
whole original buffer both fit.  Let
$\widetilde\Omega=(\id_{RC}\otimes V)\Omega$.
Its largest Schmidt probability across $Rb_R\mid Cb_C$ is at least
$e^{-b}$, hence
\[
 \Tr\rho_{\widetilde\Omega,Rb_R}^{2}\ge e^{-2b}.
\]
Pull the swap on $b_R$ back to the original doubled buffer:
\[
 W=(V^\dagger\otimes V^\dagger)F_{b_R}(V\otimes V).
\]
It is a Hermitian contraction on the two physical copies of $Q$.
The auxiliary spaces have now disappeared.  Since $V$ acts outside
$R$, the swap identity for purity gives the real positive number
\begin{equation}
 z:=\langle\Phi,W\Psi\rangle
   =\langle\widetilde\Omega^{\otimes2},
          F_RF_{b_R}\widetilde\Omega^{\otimes2}\rangle
   =\Tr\rho_{\widetilde\Omega,Rb_R}^{2}
   \ge e^{-2b}.
 \label{eq:info-reset-overlap}
\end{equation}

Write $\mathbb H$ for the sum of the original Hamiltonian on the two
replicas and $\mathbb H'=F\mathbb H F^\dagger$.  Treat the pair of
replica registers at a site as one site.  Both Hamiltonians have range
one and uniformly bounded interaction degree and strength.  Their
unique ground vectors are $\Psi,\Phi$, their common ground energy is
$\mathcal E=2E_0$, and their gaps are at least $\Delta$.  Their
difference is supported on endpoints of edges crossing $R$: terms
entirely inside $R$ merely exchange the two identical replica terms,
and all other noncrossing terms are unchanged.

Put $B_*=1+b+\log L$ and choose a fixed error exponent
$r_*=2p+10$.  For $h>0$, let
\[
 g_h(t)=(2\pi h)^{-1/2}e^{-t^2/(2h)},\qquad
 M=\int_{\mathbb R}g_h(t)e^{it\mathbb H'}W e^{-it\mathbb H}\,dt.
\]
This is a contraction.  The spectral theorem gives
\[
 M\Psi=e^{-h(\mathbb H'-\mathcal E)^2/2}W\Psi,
 \qquad
 M^\dagger\Phi=e^{-h(\mathbb H-\mathcal E)^2/2}W^\dagger\Phi.
\]
Consequently
\begin{equation}
 \norm{M\Psi-z\Phi},\quad\norm{M^\dagger\Phi-z\Psi}
       \le e^{-h\Delta^2/2}.
 \label{eq:info-gaussian-two-sided}
\end{equation}
Choose $h=A B_*$ with $A\Delta^2/2\ge10+r_*+20$.  The errors in
\eqref{eq:info-gaussian-two-sided} are then bounded by
$e^{-10b}L^{-r_*}$ up to a harmless fixed factor.  Truncate the integral
to $|t|\le T=T_0B_*$.  The resulting operator-norm error is at most
\[
 2e^{-T^2/(2h)}=2e^{-T_0^2 B_*/(2A)}.
\]
Choosing $T_0$ after $A$ with $T_0^2/(2A)\ge10+r_*+20$ gives the
same error bound.

For completeness we prove the finite-time locality estimate needed
to localize this truncated integral.  It is the usual commutator
argument for bounded finite-range interactions; see also
\cite{NachtergaeleOgataSims2006}.  For a norm-one single-site operator
$B_y$, define
\[
 c_y(S,t)=\sup_{\substack{\supp A\subseteq S\\\norm A\le1}}
                  \norm{[\tau_t(A),B_y]},
 \qquad \tau_t(A)=e^{itK}Ae^{-itK},
\]
where $K$ is either doubled Hamiltonian.  If its interaction terms are
$k_Z$, then, for $t\ge0$,
\begin{equation}
 c_y(S,t)\le c_y(S,0)+
 2\sum_{Z:Z\cap S\ne\varnothing}\norm{k_Z}
                      \int_0^t c_y(Z,v)\,dv.
 \label{eq:info-commutator-recursion}
\end{equation}
Indeed $[K,A]$ contains only interactions meeting the initial support
$S$.  Differentiate $[\tau_t(A),B_y]$ and use Jacobi's identity.
The homogeneous part is commutation by the Hermitian operator
$\sum_{Z:Z\cap S\ne\varnothing}\tau_t(k_Z)$ and is removed by a
unitary change of frame.  The remaining term for $Z$ has norm at most
$2\norm A\norm{k_Z}c_y(Z,t)$, proving
\eqref{eq:info-commutator-recursion}.  Negative times have the same
bound.

Iteration produces chains of overlapping interaction supports with
factorial time denominators.  Only chains reaching $y$ contribute an
initial commutator.  The first support has at most $C|S|$ choices,
each subsequent support has at most a fixed number of choices, and a
chain reaching distance $d_K(S,y)$ must have length at least a fixed
multiple of that distance, up to an additive constant.  Here $d_K$
is the physical site-graph distance; for either doubled Hamiltonian
it is the original nearest-neighbor distance.  Bounding the resulting
exponential series with a weight on chain length yields constants
$v,\mu,C>0$ such that
\begin{equation}
 \norm{[\tau_t(A),B_y]}
       \le C|S|\norm A\exp\bigl(v|t|-\mu d_K(S,y)\bigr).
 \label{eq:info-lr}
\end{equation}
The remainder of the iteration vanishes by its factorial denominator.
The constants depend only on the fixed local interaction bounds.

For a site set $D$, let $\mathcal P_D$ be conditional expectation
onto operators supported on $D$, obtained by averaging over independent
unitaries at all sites outside $D$.  It contracts operator norm and
preserves Hermitian operators.  Successively averaging single sites
and using \eqref{eq:info-lr} gives
\begin{equation}
 \norm{\tau_t(A)-\mathcal P_{N_\ell(S)}\tau_t(A)}
       \le C|S|L^2\norm A\,e^{v|t|-\mu\ell}.
 \label{eq:info-operator-localization}
\end{equation}
No Hilbert-space dimension appears in this averaging bound.

Let $Z_R$ be the endpoints of edges crossing $R$, and put
\[
 V(t)=e^{it\mathbb H'}e^{-it\mathbb H}.
\]
Its left generator is
\[
 \frac{d}{dt}V(t)=iG(t)V(t),\qquad
 G(t)=e^{it\mathbb H'}(\mathbb H'-\mathbb H)e^{-it\mathbb H'}.
\]
The difference $\mathbb H'-\mathbb H$ is a sum of $O(L^2)$ uniformly
bounded terms, each on at most two sites of $Z_R$.  Apply
\eqref{eq:info-operator-localization} to those terms separately and
replace $G(t)$ by
$G_\ell(t)=\mathcal P_{N_\ell(Z_R)}G(t)$.  This replacement is
Hermitian, and
\[
 \norm{G(t)-G_\ell(t)}\le CL^4e^{v|t|-\mu\ell}.
\]
The evolution $V_\ell(t)$ generated by $G_\ell(t)$, with
$V_\ell(0)=\id$, is unitary and supported on $N_\ell(Z_R)$.  Duhamel's formula for two unitary
evolutions therefore gives
\[
 \norm{V(t)-V_\ell(t)}\le CL^4|t|e^{v|t|-\mu\ell}.
\]
Localize $e^{it\mathbb H'}We^{-it\mathbb H'}$ by conditional
expectation onto $N_\ell(Q)$.  It stays a contraction, and its error
is at most $CL^4e^{v|t|-\mu\ell}$.  The factorization
\[
 e^{it\mathbb H'}We^{-it\mathbb H}
   =\bigl(e^{it\mathbb H'}We^{-it\mathbb H'}\bigr)V(t)
\]
thus has a contraction approximation supported on
\[
 S=N_\ell(Q\cup Z_R)
\]
with operator-norm error at most
$CL^4(1+T)e^{vT-\mu\ell}$ uniformly for $|t|\le T$.
If $Q$ or $Z_R$ is empty, its factor is already a scalar or the
identity and needs no localization.

Take $\ell=\lceil K_{\rm loc}B_*\rceil$.  Because
$\log L\le B_*$ and $\log(1+T)\le C+B_*$, a sufficiently large fixed
$K_{\rm loc}$, chosen after $T_0$, makes the last error at most
$Ce^{-10b}L^{-r_*}$.  Averaging the localized integrands against the
truncated Gaussian gives a square operator $M_{\rm loc}$ supported on
$S$, with $\norm{M_{\rm loc}}\le1$.  Combining the three errors, we
have
\begin{equation}
 \norm{M_{\rm loc}\Psi-z\Phi},\quad
 \norm{M_{\rm loc}^\dagger\Phi-z\Psi}
      \le C e^{-10b}L^{-r_*}.
 \label{eq:info-local-two-sided}
\end{equation}

Set $D=M_{\rm loc}/z$.  We arranged both estimates in
\eqref{eq:info-local-two-sided} so that, besides mapping $\Psi$ close
to $\Phi$, $D$ is almost isometric on $\Psi$: applying $D^\dagger$
to the first error and adding the second controls
$(D^\dagger D-\id)\Psi$.  By
\eqref{eq:info-reset-overlap} and \eqref{eq:info-local-two-sided},
\begin{align}
 \norm D&\le e^{2b},\\
 \norm{D\Psi-\Phi},\quad\norm{D^\dagger\Phi-\Psi}
      &\le C e^{-8b}L^{-r_*},\\
 \norm{(D^\dagger D-\id)\Psi}
      &\le C e^{-6b}L^{-r_*}.
 \label{eq:info-reset-polar-errors}
\end{align}
For the last estimate, write the defect as
$D^\dagger(D\Psi-\Phi)+(D^\dagger\Phi-\Psi)$.
The polar partial isometry of $D$, as a square operator on the finite
space $\mathcal H_S$, extends to a unitary $U$ on that same space.
Indeed its two support subspaces have the same dimension, as do their
orthogonal complements.  Tensor this unitary with the identity outside
$S$.  If $|D|=(D^\dagger D)^{1/2}$, then the scalar inequality
$(\sqrt{x}-1)^2\le(x-1)^2$ for $x\ge0$ gives
\[
 \norm{(U-D)\Psi}=\norm{(\id-|D|)\Psi}
       \le\norm{(\id-D^\dagger D)\Psi}.
\]
Consequently
\begin{equation}
 \norm{U\Psi-\Phi}\le C e^{-6b}L^{-r_*}\le CL^{-r_*}.
 \label{eq:info-reset-unitary}
\end{equation}
In particular, the large norm of $D$ has been explicitly paid for in
\eqref{eq:info-reset-polar-errors}; no unitary dilation on an additional
system is used.

We check the support against both cores.  Every site of $Q$ is more
than $w$ from $R_0$ and from $C_0$.  For an edge crossing $R$, write its
endpoints as $x\in R$ and $y\notin R$.  The point $y$ is more than $w$
from $R_0$, so $x$ is more than $w-1$ from $R_0$.  Since $R\cap C$ is
empty, $x\in C^c$ is more than $w$ from $C_0$, so $y$ is more than
$w-1$ from $C_0$.  Thus every site of $Z_R$ is more than $w-1$ from
both cores.  Choose $C_p$ in \eqref{eq:info-reset-width} larger than
$K_{\rm loc}+3$.  Then $S=N_\ell(Q\cup Z_R)$ is disjoint from both cores in
both replicas, including the rounding of $\ell$.

The unitary $U$ therefore leaves their joint two-copy marginal unchanged.
In $\Phi=F_R\Psi$, retaining $R_0$ and $C_0$ from one replica gives
exactly $\rho_{R_0}\otimes\rho_{C_0}$, because $R_0$ came from the
other original replica and $C_0$ did not.  Equation~\eqref{eq:info-reset-unitary}
and trace-norm contractivity under partial trace show that the original
$R_0C_0$ marginal is within trace distance $CL^{-r_*}$ of this product.
The conditional entropy continuity bound
\[
 |S(A\mid B)_\rho-S(A\mid B)_\sigma|
 \le 2\delta\log\dim\mathcal H_A
       +(1+\delta)h\left(\frac{\delta}{1+\delta}\right)
 \le C\sqrt\delta\log(e\dim\mathcal H_A),
\]
where $\delta=\tfrac12\norm{\rho-\sigma}_1$ and $h$ is binary entropy,
is valid for normalized finite-dimensional states
\cite{Winter2016}.  Applied to $I(R_0:C_0)=S(R_0)-S(R_0\mid C_0)$,
it costs at most $C\sqrt\delta\,L^2$, since $q$ is fixed.  We conclude
\[
 I_\Omega(R_0:C_0)\le CL^{2-r_*/2}=CL^{-p-3}\le L^{-p}
\]
for sufficiently large $L$.  This proves the Lemma.
\end{proof}

\subsection{Regions tapering toward finitely many points}

\begin{proposition}[Information under angular separation]
\label{prop:angular}
Fix $a>0$, $A_0\ge1$, an integer $m_0\ge0$, and $p>0$.
There is $D<\infty$ such that the following holds for all sufficiently
large $L$.  Let $\mathcal V$ be a set of at most $m_0$ marked points
in the plane, let $X,E\subseteq\Lambda$ be disjoint, and let
\[
 D\log L\le n\le A_0L,
 \qquad \diam_\infty X\le A_0n.
\]
Put $s_{\min}=D\log L$ and
$d_{\mathcal V}(x)=\min_{v\in\mathcal V}|x-v|_\infty$, with
$d_{\varnothing}=+\infty$.  Suppose that for every $x\in X$,
\begin{equation}
 \dist_\infty(x,E)
   \ge a\max\{s_{\min},\min(n,d_{\mathcal V}(x))\}.
 \label{eq:info-angular-separation}
\end{equation}
Then
\[
 I_\Omega(X:E)\le L^{-p}.
\]
The constant $D$ and the scale threshold are uniform in the marked
points, physical sets, and Hamiltonian.
\end{proposition}

\begin{proof}
An empty target or exterior is immediate.  Fix a small constant
$a'>0$, for example $a'<\min(a/100,1/100)$, and then fix
$0<d<a'/4$.  These choices depend only on the fixed geometric data.
Choose dyadically decreasing scales
\[
 s_j=2^{-j}n\quad(0\le j\le m),
 \qquad s_{\min}\le s_m<2s_{\min}.
\]
We process $X$ from farthest to nearest the marked points.  If $m\ge1$,
the first layer is $\{x\in X:d_{\mathcal V}(x)\ge n\}$, including all
greater distances.  For $1\le j<m$, the new layer is
$\{x\in X:s_j\le d_{\mathcal V}(x)<2s_j\}$.  The last layer consists
of all remaining points, namely those with
$d_{\mathcal V}(x)<2s_m$.  When $m=0$, there is just one layer, equal
to $X$.  Denote the cumulative processed set at scale $s_j$ by $A_j$.
In every case \eqref{eq:info-angular-separation} implies
\begin{equation}
 \dist_\infty(A_j,E)\ge \frac a2 s_j.
 \label{eq:info-cumulative-separation}
\end{equation}
For all nonfinal layers one may replace $a/2$ by $a$; the factor $1/2$
allows the last scale to be as large as $2s_{\min}$.

Each new layer has a cover by a bounded number of ambient squares of
side $s_j$.  Away from the first layer, it lies in the union of the
radius-$2s_j$ squares about the at most $m_0$ marked points.  For the
first or sole layer the diameter bound on $X$ gives such a cover.
The bound depends on $A_0,m_0$ but not on the shapes of the layers.

After processing scale $s=s_j$, we establish the invariant
\begin{equation}
 I_\Omega(\Lambda\cap A_j^{+a's}:\Lambda\cap E^{+a's})\le L^{-p}.
 \label{eq:info-annular-invariant}
\end{equation}
Suppose first that the preceding scale $2s$ has been processed; write
$A_{\rm old}=A_{j-1}$.  At the first scale take
$A_{\rm old}=\varnothing$ and use the same argument with an empty
initial group.  Use an ambient grid of cells of side $ds$, and let
$R$ be the physical sample of the union of all cells touching the
ambient radius-$1.5a's$ enlargement of $A_j$.  Put
\[
 C=\Lambda\cap E^{+1.5a's}.
\]
Among the selected cells, first group together those touching the
radius-$1.5a's$ enlargement of $A_{\rm old}$.  Every such cell lies in
the ambient radius-$(1.5a'+d)s$ enlargement of $A_{\rm old}$ and hence
inside its previous radius-$2a's$ enlargement.  Also
$C\subseteq\Lambda\cap E^{+2a's}$.  By the preceding invariant and data
processing, the mutual information of the entire initial group with
$C$ is at most $L^{-p}$.

Every remaining selected cell touches the radius-$1.5a's$ enlargement
of the new layer: the enlargement of a union is the union of its
enlargements.  By the bounded square cover of that layer and the fixed
ratio $d$, there are only boundedly many remaining cells.  For any
selected cell, \eqref{eq:info-cumulative-separation} and the definition
of $R,C$ give its distance from $C$ as at least
\begin{equation}
 \bigl(a/2-3a'-d\bigr)s>0.
 \label{eq:info-cell-clearance}
\end{equation}
This is a fixed positive multiple of its side $ds$.  In particular
$R$ and $C$ are disjoint.  Add the remaining cells one at a time.
The previously included physical sites are a union of whole cells,
so the chain rule and Lemma~\ref{lem:cell-cmi} apply at each addition.
For all sufficiently large $L$, $ds\ge dD\log L$ is above that
Lemma's fixed scale threshold.  Thus, after absorbing the bounded
number of additions and the initial $L^{-p}$,
\begin{equation}
 I_\Omega(R:C)\le C_b(1+\log s)^K.
 \label{eq:info-before-reset}
\end{equation}
Here $C_b,K$ are fixed independently of $D,L$ and of the current layer.

Apply Lemma~\ref{lem:reset} with
\[
 R_0=\Lambda\cap A_j^{+a's},\qquad C_0=\Lambda\cap E^{+a's}.
\]
The set $R$ contains the physical radius-$1.5a's$ enlargement of
$A_j$.  Both cores therefore have distance at least $a's/2$ from
the respective physical complements of $R,C$.  We can safely use
any strictly smaller available width $c_*s$, for instance $c_*=a'/3$.
In view of \eqref{eq:info-before-reset}, the width required by the
reset is bounded by
\begin{equation}
 w(s)\le C_3(1+\log s)^K+C_4\log L.
 \label{eq:info-reset-scale-cost}
\end{equation}
We verify uniformly over all layers that this is less than $c_*s$.
First fix all constants in the cell bound and reset, including $K$,
$C_3,C_4,c_*$.  Choose the fixed number
\[
 D\ge\frac{4C_4}{c_*}.
\]
The function $f(s)=(1+\log s)^K/s$ has logarithmic derivative
\[
 \frac{f'(s)}{f(s)}
     =\frac1s\left(\frac K{1+\log s}-1\right),
\]
so it is decreasing for $s>e^{K-1}$.  Increase the fixed threshold
$L_0$ so that, for all $L\ge L_0$,
\[
 D\log L\ge e^{\max(K,1)},\qquad
 C_3\frac{(1+\log(D\log L))^K}{D\log L}\le\frac{c_*}{2}.
\]
Such a threshold exists because
$(\log\log L)^K/\log L\to0$ for every fixed $K$.
For every scale $s\ge D\log L$ it follows that
\[
 \frac{w(s)}s
 \le C_3 f(D\log L)+\frac{C_4}{D}
 \le\frac{3c_*}{4}<c_*.
\]
This argument allows $K$ to be arbitrarily large but fixed; scales near
$e^K$ lie below the smallest processed scale once $L\ge L_0$.  No
constant or exponent is chosen as a function of $L$.

The hypotheses of Lemma~\ref{lem:reset} now hold, and its conclusion
is exactly \eqref{eq:info-annular-invariant}.  This proves the invariant
by induction.  Notice that the information cost is reset to $L^{-p}$
after each layer, so it does not accumulate over the number of layers.
At the final layer $A_m=X$.  Discarding the padding from both systems
in \eqref{eq:info-annular-invariant} gives
$I_\Omega(X:E)\le L^{-p}$, as required.
\end{proof}
